\documentclass[12pt]{article}
\usepackage{amsmath, amssymb}
\usepackage[margin=1in]{geometry}
\setlength{\parskip}{6pt}
\title{QUBO Formulation: VLSI Bisection Placement Problem}
\date{}
\begin{document}
\maketitle

\subsection*{VLSI Bisection Placement Problem}

\paragraph{Problem Statement.}
Given a set of $N$ circuit modules and a set of weighted net connections $E$ between them, the VLSI Bisection Placement Problem requires partitioning the modules into two placement regions, $A$ and $B$, such that each region contains exactly $N/2$ modules. The objective is to minimize the total weight of nets that connect a module in region $A$ to a module in region $B$.

\paragraph{Decision Variables.}
For each module $i \in \{0, 1, \dots, N-1\}$, we introduce a binary variable $x_i \in \{0, 1\}$. The value $x_i = 0$ indicates that module $i$ is assigned to region $A$, while $x_i = 1$ indicates assignment to region $B$.

\paragraph{QUBO Derivation.}
Step 1 --- Write the original objective in general notation. A net $(i,j) \in E$ with weight $w_{ij}$ crosses the bisection if and only if $x_i \neq x_j$. This condition is captured by the algebraic expression $x_i + x_j - 2x_i x_j$, which evaluates to $1$ when $x_i \neq x_j$ and $0$ otherwise. The minimization objective is therefore:
\begin{align}
    \min_{x \in \{0,1\}^N} \sum_{(i,j) \in E} w_{ij} (x_i + x_j - 2x_i x_j).
\end{align}

Step 2 --- Write each constraint in its original algebraic form. The requirement that both regions contain exactly half of the modules is expressed as:
\begin{align}
    \sum_{i=0}^{N-1} x_i = \frac{N}{2}.
\end{align}

Step 3 --- Convert the constraint to a penalty term. We enforce the constraint by adding a quadratic penalty term with weight $\lambda > 0$:
\begin{align}
    P(x) = \lambda \left( \sum_{i=0}^{N-1} x_i - \frac{N}{2} \right)^2.
\end{align}
Expanding the square and applying the idempotent property $x_i^2 = x_i$ for binary variables yields:
\begin{align}
    P(x) &= \lambda \left( \sum_{i=0}^{N-1} x_i^2 - N \sum_{i=0}^{N-1} x_i + \frac{N^2}{4} + \sum_{i \neq j} x_i x_j \right) \\
    &= \lambda \left( (1 - N) \sum_{i=0}^{N-1} x_i + \sum_{i \neq j} x_i x_j + \frac{N^2}{4} \right).
\end{align}

Step 4 --- Combine objective and penalty terms into the single QUBO Hamiltonian $H(x)$:
\begin{align}
    H(x) = \sum_{(i,j) \in E} w_{ij} (x_i + x_j - 2x_i x_j) + \lambda \left( (1 - N) \sum_{i=0}^{N-1} x_i + \sum_{i \neq j} x_i x_j + \frac{N^2}{4} \right).
\end{align}

Step 5 --- Read off the Q matrix entries and offset $c$ from $H$ in general closed form. Collecting coefficients for each monomial $x_i x_j$ (with $i \leq j$) gives:
\begin{align}
    Q_{i,i} &= \sum_{j : (i,j) \in E} w_{ij} + \lambda(1 - N), \quad \forall i \in \{0, \dots, N-1\}, \\
    Q_{i,j} &= -2w_{ij} + 2\lambda, \quad \forall i < j \text{ such that } (i,j) \in E, \\
    Q_{i,j} &= 2\lambda, \quad \forall i < j \text{ such that } (i,j) \notin E, \\
    c &= \lambda \frac{N^2}{4}.
\end{align}

\paragraph{Penalty Weight Justification.}
The penalty coefficient $\lambda$ must be chosen sufficiently large to ensure that any violation of the balance constraint incurs a higher energy cost than the maximum possible reduction in the cut weight. Specifically, let $W_{\max} = \max_{(i,j) \in E} w_{ij}$ denote the maximum net weight. Choosing $\lambda > W_{\max}$ guarantees that the penalty for deviating from the required partition size strictly dominates any objective improvement achievable by an unbalanced assignment.

\paragraph{Correctness Argument.}
Minimizing $H(x)$ over $\{0,1\}^N$ recovers the optimal bisection because feasible solutions (those satisfying $\sum_i x_i = N/2$) incur zero penalty and their energy equals the cut weight. Infeasible solutions incur a penalty of at least $\lambda$. Since $\lambda > W_{\max}$, the energy of any infeasible solution exceeds the energy of the best feasible solution by a margin that cannot be offset by reducing the cut weight. Consequently, the global minimum of $H(x)$ must correspond to a feasible assignment that minimizes the total cut weight.

\paragraph{Illustrative Example.}
Consider a minimal instance with $N=2$ modules and a single net connecting module $0$ and module $1$ with weight $w_{01}$. The constraint requires exactly one module in region $B$, i.e., $x_0 + x_1 = 1$. Substituting $N=2$ into the general formulas yields the concrete penalty expansion:
\begin{align}
    P(x) = \lambda \left( (1-2)(x_0 + x_1) + 2x_0 x_1 + 1 \right) = \lambda \left( 1 - x_0 - x_1 + 2x_0 x_1 \right).
\end{align}
Combining this with the objective $w_{01}(x_0 + x_1 - 2x_0 x_1)$ gives the full Hamiltonian:
\begin{align}
    H(x) = (w_{01} + \lambda)(x_0 + x_1) - 2(w_{01} + \lambda)x_0 x_1 + \lambda.
\end{align}
The corresponding QUBO parameters are:
\begin{align}
    Q_{0,0} &= w_{01} + \lambda, & Q_{1,1} &= w_{01} + \lambda, \\
    Q_{0,1} &= -2(w_{01} + \lambda), & c &= \lambda.
\end{align}
Evaluating $H(x)$ over all four binary assignments reveals that the configurations $(x_0, x_1) = (0,1)$ and $(1,0)$ yield energy $w_{01} + \lambda - 2(w_{01} + \lambda) + \lambda = -w_{01} + \lambda$, while the balanced configurations $(0,0)$ and $(1,1)$ yield energy $w_{01} + \lambda$. Since $w_{01} > 0$, the minimum energy is $-w_{01} + \lambda$, correctly identifying the optimal bisection where the single net is cut, with minimum cut weight $w_{01}$.

\end{document}